% Informe tecnico - CODEFEST AD ASTRA 2026, Etapa 1
% Compilar:  xelatex informe_tecnico.tex
\documentclass[10pt,a4paper]{article}

\usepackage{fontspec}
\usepackage[spanish,es-noquoting,es-nolists]{babel}
\usepackage[margin=2.1cm]{geometry}
\usepackage{booktabs}
\usepackage{amsmath}
\usepackage{xcolor}
\usepackage{titlesec}
\usepackage{enumitem}
\usepackage{hyperref}
\usepackage{fancyhdr}
\usepackage{tabularx}

\definecolor{azul}{RGB}{20,60,120}
\definecolor{rojo}{RGB}{170,30,45}
\definecolor{gris}{RGB}{90,90,90}

\hypersetup{colorlinks=true, linkcolor=azul, urlcolor=azul}

\titleformat{\section}{\large\bfseries\color{azul}}{\thesection.}{0.5em}{}
\titleformat{\subsection}{\normalsize\bfseries\color{rojo}}{\thesubsection}{0.5em}{}
\titlespacing*{\section}{0pt}{1.1ex}{0.6ex}
\titlespacing*{\subsection}{0pt}{0.9ex}{0.4ex}

\setlist[itemize]{leftmargin=1.2em, itemsep=0.15ex, topsep=0.3ex}
\setlist[enumerate]{leftmargin=1.4em, itemsep=0.15ex, topsep=0.3ex}
\setlength{\parskip}{0.35ex}

\pagestyle{fancy}
\fancyhf{}
\fancyhead[L]{\footnotesize\color{gris}CODEFEST · AD ASTRA 2026}
\fancyhead[R]{\footnotesize\color{gris}Etapa 1 — Base de Conocimiento}
\fancyfoot[C]{\footnotesize\thepage}
\renewcommand{\headrulewidth}{0.4pt}

\begin{document}

\begin{center}
{\LARGE\bfseries\color{azul} Base de Conocimiento Vectorial}\\[0.3ex]
{\large Informe Técnico — Etapa 1}\\[0.5ex]
{\footnotesize\color{gris} CODEFEST AD ASTRA 2026 · Fuerza Aeroespacial Colombiana · Universidad de los Andes}
\end{center}
\vspace{0.4ex}
\hrule
\vspace{1ex}

\section{Resumen}

Construimos un sistema de \textbf{recuperación densa multilingüe} que indexa un corpus multiformato
(ES/EN/PT) sobre los tres fenómenos del reto y responde consultas en lenguaje natural devolviendo
los 3 documentos y los 10 fragmentos más relevantes.

El diseño parte de una lectura precisa de las reglas de evaluación (Sección 10.2.1 de la
especificación): \textbf{la relevancia de un fragmento se juzga sobre su contenido textual}
(\texttt{text}) \textbf{y la de un documento sobre el campo} \texttt{fuente}, no sobre los
identificadores internos que asigna cada equipo. De ahí nuestra prioridad de ingeniería: la
\emph{fidelidad de la extracción y la limpieza del texto} determinan el techo de las métricas,
por encima de cualquier ajuste fino del recuperador.

\begin{center}\small\ttfamily
Extracción $\rightarrow$ Limpieza $\rightarrow$ Chunking (2 niveles) $\rightarrow$ Encoder(s) $\rightarrow$ FAISS\\
$\rightarrow$ Fusión RRF $\rightarrow$ [Reranking] $\rightarrow$ Agregación chunk$\rightarrow$doc $\rightarrow$ resultados.jsonl\\
(+ Grafo de conocimiento)
\end{center}

\section{Preprocesamiento: extracción y limpieza}

\subsection{Extracción por formato}

\begin{center}\small
\begin{tabularx}{\textwidth}{@{}l l X@{}}
\toprule
\textbf{Formato} & \textbf{Herramienta} & \textbf{Motivo} \\
\midrule
PDF & Docling (MIT) + PyMuPDF & Docling reconstruye layout, tablas y orden de lectura; PyMuPDF garantiza robustez si Docling falla \\
HTML & trafilatura & Elimina \emph{boilerplate} conservando el cuerpo del artículo \\
JSON & selección de campos & Concatena \texttt{title}/\texttt{body\_text} en orden; \texttt{url}, \texttt{date} van a metadata \\
CSV/XLSX & pandas + openpyxl & Cada fila serializada como pares \texttt{columna: valor} para preservar contexto \\
Imágenes & OCR multilingüe & Recupera texto de infografías y figuras \\
PBF & recorrido de capas & Atributos como pares clave-valor; una versión por elemento (evita duplicar por zoom) \\
\bottomrule
\end{tabularx}
\end{center}

La elección de Docling se validó empíricamente: sobre un PDF institucional real de 24 páginas
extrajo \textbf{13\,\% más contenido} que PyMuPDF y fue el único que recuperó una tabla como
estructura tabular en lugar de texto aplanado.

\subsection{Limpieza y normalización}

Esta etapa resultó la de \textbf{mayor impacto medido} del preprocesamiento. Al analizar la
extracción de un PDF real detectamos tres patrones de ruido que degradaban el índice:

\begin{enumerate}
\item \textbf{Índices con puntos de relleno} (\texttt{Introducción . . . . 3}): el segmentador
interpretaba cada punto como fin de oración.
\item \textbf{Encabezados y pies repetidos} en cada página, que generaban fragmentos sin valor
informativo y diluían la señal semántica.
\item \textbf{Guiones de corte de línea} (\texttt{conges-\textbackslash n tión}), que parten
palabras y rompen el emparejamiento léxico y semántico.
\end{enumerate}

Nuestra limpieza normaliza a NFC, elimina caracteres de control, descarta líneas de índice,
detecta \emph{boilerplate} por repetición y \textbf{re-une palabras cortadas por guion},
distinguiéndolas de compuestos legítimos (\texttt{space-debris} no se altera).

\textbf{Efecto medido:} en los primeros 3\,000 caracteres del documento de prueba, las falsas
oraciones detectadas pasaron de \textbf{279 a 59}, y desaparecieron los fragmentos compuestos
únicamente por el índice del documento.

\section{Estrategia de chunking}

\subsection{Chunking de dos niveles}

La especificación impone dos restricciones que operan en momentos distintos: la
\textbf{completitud lingüística} (Sección 3.3), que prohíbe fragmentos con oraciones incompletas,
y el \textbf{límite de 250 palabras} por fragmento de salida (Sección 9.2). El tamaño óptimo para
\emph{codificar} no coincide con el máximo de \emph{presentación}, así que los separamos:

\begin{itemize}
\item \textbf{Nivel 1 — chunk de índice.} Agrupa oraciones completas hasta un máximo configurable
de tokens (256–384, dentro del límite del encoder), con solapamiento de una oración entre
fragmentos consecutivos. Es la unidad que se codifica y almacena en FAISS.
\item \textbf{Nivel 2 — sub-fragmento de salida.} Al construir la respuesta, un chunk que exceda
las 250 palabras se divide respetando de nuevo los límites oracionales. Los sub-fragmentos
conservan el \texttt{chunk\_id} original, que cumple función de \textbf{trazabilidad} hacia la
evidencia indexada.
\end{itemize}

\subsection{Justificación}

\begin{itemize}
\item \emph{Frente a tamaño fijo de tokens:} cortar por conteo parte oraciones y viola el
requisito obligatorio de la Sección 3.3.
\item \emph{Frente a chunking por párrafo puro:} los párrafos de documentos técnicos van de una
línea a más de una página; produce fragmentos desiguales que pueden superar la ventana del encoder.
\item \emph{A favor del solapamiento:} una idea en la frontera entre dos fragmentos quedaría
representada a medias en ambos; solapar una oración lo mitiga con bajo coste.
\end{itemize}

La segmentación oracional usa \textbf{pysbd} (multilingüe), con un segmentador de respaldo basado
en reglas para que el sistema opere aunque falte la dependencia.

\subsection{Metadata por fragmento}

Cada chunk almacena los ocho campos obligatorios de la Tabla 1 más campos opcionales
(\texttt{idioma}, \texttt{titulo}, \texttt{fecha}). El campo \texttt{fuente} se normaliza siempre
con separadores \texttt{/}: durante el desarrollo detectamos que en Windows se guardaba con
\texttt{\textbackslash}, lo que \textbf{anulaba el emparejamiento a nivel documento} y llevaba el
F1@3 a cero pese a una recuperación correcta.

\section{Codificación semántica: selección de encoder}

\subsection{Encoder principal: BAAI/bge-m3}

\begin{center}\small
\begin{tabularx}{\textwidth}{@{}l X@{}}
\toprule
\textbf{Criterio (Sección 4.3)} & \textbf{Cumplimiento} \\
\midrule
Soporte multilingüe & Nativo en ES/EN/PT, entrenado para recuperación cross-lingual \\
Dimensionalidad & 1024 — equilibrio entre expresividad y coste \\
Longitud de entrada & Hasta 8\,192 tokens, muy por encima de los 512 habituales \\
Benchmarks & Estado del arte en MIRACL/MTEB multilingüe para recuperación densa \\
Licencia & MIT \\
Eficiencia & Ejecuta en una GPU de 8\,GB en fp16 \\
\bottomrule
\end{tabularx}
\end{center}

El criterio decisivo es el \textbf{cross-lingual}: el conjunto de evaluación reparte las consultas
entre español, inglés y portugués, y una consulta en un idioma debe recuperar documentos en los
otros dos. Un encoder monolingüe fallaría en dos tercios de las consultas.

\textbf{Validación empírica.} BGE-M3 obtuvo NDCG@10 de \textbf{0,990} frente a \textbf{0,966} de
\texttt{paraphrase-multi\-lingual-MiniLM-L12-v2}, con la diferencia concentrada exactamente donde
se predijo: la consulta en portugués pasó de 0,877 a \textbf{1,000}.

\subsection{Restricción de arquitectura}

Todos los modelos empleados en indexación y recuperación son \textbf{encoders} (familia BERT). No
se utiliza ningún modelo generativo (decoder) en la construcción del índice ni en la recuperación,
conforme a las Secciones 4.2 y 8.3.

\section{Índice vectorial FAISS}

Elegimos \textbf{\texttt{IndexFlatIP} con vectores normalizados a norma unitaria}, que equivale
exactamente a similitud coseno (ecuación 4 de la especificación).

\textbf{Justificación:} para el volumen de este reto la búsqueda exhaustiva es holgadamente viable
y ofrece resultados \textbf{exactos}. Los índices aproximados (\texttt{IVFFlat}, \texttt{HNSW})
sacrifican exactitud por velocidad, un intercambio que no compensa cuando la métrica evaluada es la
calidad del ranking y el tiempo de respuesta no se puntúa. Si el corpus creciera un orden de
magnitud, migrar a \texttt{IndexHNSWFlat} es un cambio de una línea en la configuración.

FAISS almacena únicamente vectores e identificadores internos. La metadata vive en
\texttt{metadata.jsonl}, donde \textbf{la línea $i$ corresponde al identificador interno $i$} que
FAISS asigna al indexar, tal como exige el formato de entrega. La persistencia usa las funciones
nativas \texttt{write\_index}/\texttt{read\_index}.

\section{Módulo de recuperación}

\subsection{Fusión de múltiples índices}

Implementamos las tres estrategias de la Sección 8.4 y adoptamos \textbf{Reciprocal Rank Fusion}
por defecto:
\[ s_{\mathrm{RRF}}(c) = \sum_{j=1}^{m} \frac{1}{k_0 + r_j(c)}, \qquad k_0 = 60 \]
RRF opera sobre \textbf{posiciones} y no sobre puntuaciones, por lo que es robusto frente a las
distintas escalas de similitud de encoders diferentes —un problema real, ya que CombSUM permitiría
que el encoder con puntuaciones sistemáticamente más altas domine la fusión.

\subsection{Agregación de fragmentos a documentos}

Para el nivel documento (F1@3) agrupamos los chunks por \texttt{doc\_id} y calculamos una
puntuación agregada. Implementamos tres estrategias —máximo (\emph{max pooling}), suma y media
ponderada— comparadas empíricamente en la Sección 8.

\subsection{Reranking con cross-encoder}

Incluimos un paso opcional de reordenamiento con \texttt{BAAI/bge-reranker-v2-m3}, un
\textbf{cross-encoder} que puntúa conjuntamente el par (consulta, fragmento).

\textbf{Consideración sobre el reglamento.} La Sección 8.3 prohíbe modelos \textbf{generativos}
(decoders) en la recuperación. Un cross-encoder es arquitectura \textbf{encoder} (familia BERT) y
produce una puntuación escalar de relevancia; no genera texto. Entendemos que su uso es admisible,
pero por prudencia el paso es \textbf{desactivable por configuración} y, ante una interpretación
restrictiva, el sistema opera con fusión RRF pura sin modificar una línea de código.

\subsection{Cumplimiento del formato de salida}

Un validador estricto verifica cada objeto antes de escribir \texttt{resultados.jsonl}: exactamente
3 documentos, exactamente 10 fragmentos, ningún fragmento por encima de 250 palabras y orden
\texttt{q001}–\texttt{q050}. La verificación es automática y forma parte del pipeline de generación.

\section{Grafo de conocimiento (componente bonus)}

\textbf{Construcción.} (i) \emph{Reconocimiento de entidades} con \texttt{GLiNER} multilingüe
(MIT), modelo \emph{zero-shot} que extrae tipos definidos en configuración sin reentrenamiento;
verificamos su comportamiento cross-lingual: sobre texto en español identificó \emph{Estados
Unidos} (país), \emph{sistemas de armas autónomas} (tecnología) y \emph{Convenio de Ginebra}
(evento). (ii) \emph{Extracción de relaciones} por co-ocurrencia dentro del mismo fragmento, con
peso acumulado; \textbf{no se emplea ningún modelo generativo}. (iii) \emph{Persistencia} en grafo
dirigido NetworkX exportado a \texttt{grafo.graphml}, donde cada nodo guarda los \texttt{chunk\_id}
en que aparece la entidad y cada arista referencia el \texttt{doc\_id} y \texttt{chunk\_id} que la
respalda, garantizando \textbf{trazabilidad} de toda relación hasta su evidencia textual.

\textbf{Integración con la recuperación.} El grafo actúa como \textbf{índice adicional} dentro del
esquema de fusión (Sección 8.5): se extraen las entidades de la consulta con el mismo NER, se
recuperan los chunks vinculados a esas entidades y a sus \textbf{vecinos de primer orden}, se
puntúan por evidencia acumulada y el ranking resultante se fusiona con los vectoriales mediante
RRF. Aporta una señal que los embeddings sólo capturan de forma implícita: \textbf{relaciones
explícitas entre entidades}.

\section{Metodología de evaluación y resultados}

\subsection{Optimizar sin ground truth}

El conjunto de evaluación y sus juicios de relevancia no son públicos durante el reto. Ajustar
hiperparámetros «a ojo» equivale a no ajustarlos. Construimos un \textbf{arnés de evaluación
interno}: (i) consultas propias con relevancia conocida, equilibradas entre fenómenos e idiomas;
(ii) una implementación de \textbf{NDCG@10 y F1@3 idéntica a la especificación} (ecuaciones 8–14),
verificada con pruebas unitarias que reproducen el ejemplo de Conteo de Borda de la Tabla 3;
(iii) un barrido que rankea cada configuración con \textbf{Conteo de Borda}, el mismo criterio del
leaderboard oficial, optimizando así el equilibrio entre ambas métricas y no una sola.

Este arnés detectó dos defectos que habrían costado puntos y que ninguna inspección visual habría
revelado: el separador de rutas en \texttt{fuente} y un error en el ranking ideal del NDCG que
producía valores superiores a 1.

\subsection{Corpus de validación previo}

Antes de disponer del corpus oficial reunimos \textbf{151 documentos públicos auténticos}
(artículos de arXiv, \emph{ESA Space Environment Report}, informe IADC de UNOOSA, \emph{Panorama
Social} de la CEPAL) y construimos 34 consultas con juicios de relevancia derivados de una señal
externa e independiente: el propio buscador de arXiv. Si una consulta devolvió un artículo, ese
artículo es relevante para la consulta en lenguaje natural equivalente, y su posición da el grado.
Es supervisión débil, pero ajena a nuestro sistema, lo que evita el sesgo circular de evaluarnos
con juicios propios.

\subsection{El corpus oficial}

El corpus entregado consta de \textbf{1\,826 documentos} (2,93\,GB) de \textbf{20 observatorios}.
Difería de nuestras previsiones en tres puntos que obligaron a adaptar el sistema:

\begin{itemize}
\item \textbf{El formato dominante es JSON} (964 archivos, 52\,\%), no PDF (759). Los esquemas son
heterogéneos —artículos con \texttt{body\_text}, páginas con \texttt{sections}, catálogos de
fichas—, así que reescribimos el extractor para recorrer cualquier estructura y recoger todo el
texto útil descartando metadata técnica, en lugar de buscar un conjunto fijo de claves.
\item \textbf{ADL entrega un inventario oficial} con un \texttt{DOC\_ID} único por archivo. Usarlo
elimina la mayor incertidumbre del reto: la evaluación a nivel documento empareja por
\texttt{fuente}, y así nos apoyamos en los identificadores del organizador en vez de inventarlos.
\item \textbf{51 PDFs no tienen capa de texto.} Son escaneados y la extracción devolvía cadena
vacía \emph{sin error visible}. 47 son informes de Alertas Tempranas (Defensoría del Pueblo), el
observatorio más relevante para 18 de las 50 consultas. Se recuperaron mediante OCR.
\end{itemize}

\textbf{Resultado de la indexación:} 82\,673 fragmentos de 1\,688 documentos, metadata alineada
1:1 con el índice FAISS, índice disperso con 59\,673 tokens únicos, y \texttt{resultados.jsonl}
con \textbf{cero errores de esquema}.

\subsection{Dos hallazgos que solo aparecen con el corpus real}

\textbf{Contaminación por el archivo de consultas.} El paquete de ADL incluye, junto al corpus, el
PDF con las 50 consultas de evaluación. Al indexarlo como un documento más, emparejaba de forma
trivial —contiene el texto literal de cada consulta— y aparecía en el top-3 de \textbf{20 de las 50
consultas}, gastando uno de los tres únicos huecos de documento. La regla que lo resuelve es
\textbf{el inventario oficial define el corpus}: todo archivo ausente de él se omite. Contaminación
medida: 20/50 $\rightarrow$ 0/50.

\textbf{Concentración de fragmentos.} Tres volcados bibliográficos en CSV generaban 78\,000
fragmentos entre los tres —el 51\,\% del índice— sobre biomedicina ajena a las consultas. Un
límite configurable por documento acota su peso sin excluirlos del corpus.

\subsection{Selección del motor de OCR}

Medimos dos motores sobre la misma página escaneada. PaddleOCR requiere desactivar oneDNN para no
fallar en Windows, lo que elimina sus kernels optimizados: \textbf{115\,s por página}, es decir
20 horas para las 648 páginas escaneadas del corpus. EasyOCR, que se apoya en la misma pila de
GPU ya empleada por los encoders, baja a \textbf{6,5\,s por página} (48 minutos en total) con
calidad equivalente. Ambos son modelos de detección y reconocimiento clásicos, no generativos.

\subsection{Resultados}

Los 50 documentos se procesaron \textbf{sin una sola omisión}, produciendo \textbf{2\,601
fragmentos}. El almacén de metadata pasó la validación de esquema con \textbf{cero errores} y
mantiene la correspondencia exacta línea $\leftrightarrow$ identificador interno de FAISS.

\begin{center}\small
\begin{tabular}{@{}lccc@{}}
\toprule
\textbf{Configuración} & \textbf{NDCG@10} & \textbf{F1@3} & \textbf{Borda} \\
\midrule
\textbf{max pooling + reranking} & \textbf{0,4246} & \textbf{0,5000} & \textbf{8} \\
media ponderada & 0,3625 & 0,5000 & 6 \\
max pooling & 0,3625 & 0,5000 & 6 \\
media ponderada + reranking & 0,4246 & 0,4697 & 5 \\
suma & 0,3625 & 0,4697 & 3 \\
\bottomrule
\end{tabular}
\end{center}

\textbf{Hallazgo 1 — el reranking es la palanca más rentable.} El cross-encoder eleva el NDCG@10 de
0,3625 a \textbf{0,4246}, una mejora relativa del \textbf{17\,\%}, sin degradar el F1@3. Es la
mayor ganancia individual medida en todo el sistema.

\textbf{Hallazgo 2 — los corpus de juguete engañan.} Sobre un conjunto sintético reducido, las
mismas variantes indicaban que el reranking \textbf{no aportaba nada} e incluso perjudicaba, porque
con pocos documentos el recuperador denso ya saturaba las métricas (NDCG $\approx$ 0,95 en toda
variante). De haber confiado en ese conjunto habríamos descartado el componente más valioso.

\textbf{Hallazgo 3 — la agregación por suma perjudica el nivel documento.} En ambos corpus, sumar
las puntuaciones de todos los fragmentos de un documento reduce el F1@3: favorece documentos largos
con muchos fragmentos mediocres sobre documentos concisos con un fragmento excelente. Adoptamos
\textbf{max pooling}.

\textbf{Interpretación de las cifras absolutas.} Los juicios de este conjunto son
\textbf{incompletos por construcción}: cada consulta tiene unos cinco documentos etiquetados, pero
el corpus contiene otros genuinamente relacionados que no lo están; cuando el sistema los recupera
—acertadamente— computan como error y deprimen el NDCG. Estas cifras son válidas para
\textbf{comparar configuraciones entre sí}, que es el uso que les damos, pero no son una estimación
del rendimiento absoluto sobre el conjunto oficial, cuyos juicios sí son exhaustivos.

\section{Reproducibilidad}

El script \texttt{generador.py} reconstruye \texttt{resultados.jsonl} a partir del índice
persistido y del archivo de consultas. El determinismo se asegura fijando las semillas de
\texttt{random}, NumPy y PyTorch, activando algoritmos deterministas y escribiendo el JSON Lines
con separadores fijos y codificación UTF-8 explícita. Toda la configuración —chunking, encoders,
tipo de índice, fusión, agregación— reside en un único \texttt{config.yaml}, de modo que cualquier
resultado reportado es rastreable a una configuración concreta. El proyecto cuenta con
\textbf{42 pruebas automatizadas} que cubren las métricas, el validador de esquema, el chunker, las
estrategias de fusión, la limpieza y el recuperador basado en grafo.

\section{Conclusiones}

La decisión de diseño más rentable no fue la elección del modelo, sino \textbf{leer la regla de
emparejamiento y orientar el esfuerzo a la fidelidad del texto}. La limpieza del corpus produjo la
mejora más barata; el encoder correcto (BGE-M3) resolvió el requisito cross-lingual; el reranking
aportó la mayor ganancia medida (+17\,\% de NDCG@10); y el arnés de evaluación interno convirtió
decisiones de intuición en decisiones medidas.

La lección metodológica que atraviesa el trabajo es que \textbf{medir en condiciones
representativas cambia las conclusiones}. Dos hallazgos —el valor del reranking y el perjuicio de
la agregación por suma— sólo emergieron al evaluar sobre documentos reales, largos y multilingües.
Y dos defectos que habrían costado puntos fueron detectados por las pruebas y el arnés, no por
inspección visual.

Los componentes de mayor riesgo regulatorio quedan aislados tras interruptores de configuración,
de modo que el sistema puede ajustarse a la interpretación del jurado sin tocar código. Nuestra
recomendación, sujeta a esa consulta, es mantener el reranking activo: es la diferencia medida más
grande entre una entrada correcta y una competitiva.

\end{document}
